%  LaTeX support: latex@mdpi.com 
%  For support, please attach all files needed for compiling as well as the log file, and specify your operating system, LaTeX version, and LaTeX editor.

%=================================================================
\documentclass[geomatics,article,submit,pdftex,moreauthors]{Definitions/mdpi} 

%=================================================================
%----------
% journal

% MDPI internal commands
\firstpage{1} 
\makeatletter 
\setcounter{page}{\@firstpage} 
\makeatother
\pubvolume{1}
\issuenum{1}
\articlenumber{0}
\pubyear{2024}
\copyrightyear{2024}
%\externaleditor{Academic Editor: Firstname Lastname}
\datereceived{09 December 2024} 
\daterevised{04 January 2025} 
\dateaccepted{} 
\datepublished{} 
\hreflink{https://doi.org/} % If needed use \linebreak
%\doinum{}

%=================================================================
% Add packages and commands here. The following packages are loaded in our class file: fontenc, inputenc, calc, indentfirst, fancyhdr, graphicx, epstopdf, lastpage, ifthen, lineno, float, amsmath, setspace, enumitem, mathpazo, booktabs, titlesec, etoolbox, tabto, xcolor, soul, multirow, microtype, tikz, totcount, changepage, attrib, upgreek, cleveref, amsthm, hyphenat, natbib, hyperref, footmisc, url, geometry, newfloat, caption

\graphicspath{{figures/}} % path to folder with figures
\usepackage{subcaption}
\usepackage{listings}
\usepackage{xcolor}
\usepackage{gensymb} % for \textdegree
\usepackage{tabularx}
\newcolumntype{Y}{>{\centering\arraybackslash}X}
\usepackage{amsmath,mathtools,amsthm}
	\setcounter{MaxMatrixCols}{15}
\usepackage{array}
\usepackage{amssymb}
\usepackage{amsfonts}
\usepackage{float}
\usepackage{chngcntr}
\counterwithin*{equation}{section}
\usepackage[super]{nth}
\usepackage{longtable}
\usepackage{tabularx}
\usepackage{multirow}
\usepackage{adjustbox}

\definecolor{deepblue}{rgb}{0,0,0.5}
\definecolor{deepred}{rgb}{0.6,0,0}
\definecolor{deepmagenta}{RGB}{195,12,214}
\definecolor{deepgreen}{RGB}{29,122,30}
\definecolor{mintgreen}{RGB}{192,249,192}
\definecolor{codepurple}{RGB}{204, 0, 153}
\definecolor{LightCyan}{rgb}{0.88,1,1}
\definecolor{GainsboroPurple}{RGB}{247,176,247}
\definecolor{LightSlateGrey}{RGB}{124,118,118}
%    
\lstdefinestyle{mystyle}{
    commentstyle=\color{LightSlateGrey},
    keywordstyle=\color{deepgreen}\bfseries,
    emphstyle=\color{deepred},
    numberstyle=\tiny\color{deepblue},
    stringstyle=\color{codepurple},
    basicstyle=\ttfamily\scriptsize,
    breakatwhitespace=false,         
    breaklines=true,                 
    captionpos=b,                    
    keepspaces=true,                 
    numbers=left,                    
    numbersep=5pt,                  
    showspaces=false,                
    showstringspaces=false,
    showtabs=false,                  
    tabsize=2
}
\lstset{style=mystyle}

\usepackage{subcaption}
\usepackage{listings}
\usepackage{xcolor}
\usepackage{gensymb} % for \textdegree
\usepackage{amsmath,mathtools,amsthm}
	\setcounter{MaxMatrixCols}{15}
\usepackage{amssymb}
\usepackage{amsfonts}
\usepackage{float}
\usepackage{chngcntr}
\counterwithin*{equation}{section}
\usepackage[super]{nth}
\usepackage{longtable}
\usepackage{multirow}
\usepackage{adjustbox}
\usepackage{caption}
\DeclareCaptionType{equ}[][]

%=================================================================
% Full title of the paper (Capitalized)
\Title{Improving Bimonthly Landscape Monitoring in Morocco, North Africa, by Integrating Machine Learning with GRASS GIS}

% MDPI internal command: Title for citation in the left column
\TitleCitation{Improving Bimonthly Landscape Monitoring in Morocco, North Africa, by Integrating Machine Learning with GRASS GIS}

% Author Orchid ID: enter ID or remove command
\newcommand{\orcidauthorA}{0000-0002-5759-1089} % Polina Add \orcidA{} behind the author's name

% Authors, for the paper (add full first names)
\Author{Polina Lemenkova $^{1,2}$*\orcidA{}}

%\longauthorlist{yes}

% MDPI internal command: Authors, for metadata in PDF
\AuthorNames{Polina Lemenkova}

% MDPI internal command: Authors, for citation in the left column
\AuthorCitation{Lemenkova, P.}
% If this is a Chicago style journal: Lastname, Firstname, Firstname Lastname, and Firstname Lastname.

% Affiliations / Addresses (Add [1] after \address if there is only one affiliation.)
\address{%
$^{1}$ \quad Alma Mater Studiorum -- Università di Bologna, Department of Biological, Geological and Environmental Sciences, Via Irnerio 42, Bologna, Emilia-Romagna, Italy, IT-40126. ORCID- 0000-0002-5759-1089, (P.L.).

$^{2}$ \quad Libera Università di Bolzano, Faculty of Agricultural, Environmental and Food Sciences, Piazza Università 5, IT-39100, Bolzano, Trentino-Alto Adige | S\"udtirol, Italy.
}

% Contact information of the corresponding author
\corres{Correspondence: polina.lemenkova2@unibo.it ; polina.lemenkova@unibz.it . Tel.: +39-3446928732 (P.L.)}

% Current address and/or shared authorship
%\firstnote{Current address: Affiliation 3.} 

% Abstract (Do not insert blank lines, i.e. \\) 
\abstract{This article presents the \textcolor{blue}{application of novel} cartographic methods of vegetation mapping  \textcolor{blue}{with a case study of} Rif Mountains, northern Morocco. \textcolor{blue}{The study area in notable for} varied geomorphology and diverse landscapes. \textcolor{blue}{The methodology includes ML}  modules of GRASS GIS 'r.learn.train', 'r.learn.predict' and 'r.random' \textcolor{blue}{with algorithms of supervised classification implemented from} the Scikit-Learn libraries of Python. This approach provides a platform for \textcolor{blue}{processing} spatiotemporal data \textcolor{blue}{ and satellite image analysis}. The objective is to determine the robustness of \textcolor{blue}{the “DecisionTreeClassifier” and “ExtraTreesClassifier” classification algorithms}. \textcolor{blue}{The time series} of satellite images \textcolor{blue}{covering} northern Morocco \textcolor{blue}{consists of six} Landsat \textcolor{blue}{scenes for} 2023 with a bimonthly time interval. Land cover maps are produced \textcolor{blue}{based on the} processed, classified and analysed images. \textcolor{blue}{The results demonstrated} seasonal changes in vegetation and land cover types. \textcolor{blue}{The validation was performed using land cover dataset from Food and Agriculture Organization (FAO)}. This study contributes to the \textcolor{blue}{environmental monitoring in North Africa using ML algorithms of satellite image processing. Using RS data combined with the powerful functionality of GRASS GIS and FAO-derived dataset, the topographic variability, moderate-scale habitat heterogeneity and bimonthly distribution of land cover types of northern Morocco in 2023 have been assessed for the first time.}}

% Keywords
\keyword{Machine Learning; Remote Sensing; Image Processing; Satellite Image; Mapping; Africa; Image Analysis; Automation; Geography; Earth Sciences} 

% The fields PACS, MSC, and JEL may be left empty or commented out if not applicable
%\PACS{J0101}
%\MSC{}
%\JEL{}

\begin{document}

%----------- section ----------->
\section{Introduction}

The study of the relationship between geomorphology and vegetation has become a major issue for monitoring territories using Earth observation \hl{(EO)} data \cite{Corenblit,Muir,Pu,Cary01091990}. \hl{Time series of such} geospatial data are \hl{valuable for evaluation of environmental changes based on their} classification \cite{rs15225374,w17010068,app15010338}. \hl{Nevertheless, the issue of environmental monitoring} using \hl{Remote Sensing (RS}) data consists in receiving new data faster than we are able to describe, interpret and process them. This is especially true in the \hl{satellite time series, where is the need to classify, evaluate and map habitats in operative regime} \cite{rs15225413,rs15225408,rs15225415}. In particular, geospatial data is used in the areas of agricultural land management and climate change \cite{Ahmadian03052016,Lemenkova202412,IMTIAZ2024109172}. \hl{Operative mapping of} areas with diversified geomorphology is based on \hl{time series of the }satellite data that allow\hl{ }regular \hl{visualization for environmental monitoring} \cite{JOHANSSON01011978,rs9121306}.  

In Morocco, the \hl{benefits from the use} of satellite images for landscape analysis \hl{are} explained by \hl{a} very diverse landscape structure \cite{Sahrane13122022,Marco04032021,Aouraghe01032010}. Since the geography of Morocco extends from the Atlantic Ocean to mountainous areas and the Sahara desert, \hl{this} country \hl{is notable for the diversified }geomorphological framework and varied topography. \hl{Moreover}, the geological and tectonic development of the Atlas Mountains and their surroundings presents a complex pattern of geological units \cite{Benvenuti31122024,Ouallali03072020,Khalloufi01032010}. This also determines \hl{the} distribution \hl{of major vegetation types,} and controls land use patterns in the mountainous regions of \hl{the country} \cite{ALIELOMAIRI2024101321}. 

In this context, land cover maps \hl{that visualise regional} environmental \hl{setting} play a vital role in mapping\hl{ }Morocco. \hl{Because RS data} can be interpreted in terms of landscape structural units, comparing \hl{several }images \hl{in a time series }is useful for detecting changes\hl{. In turn, change detection of land cover types enables to reveal major steps of }landscape dynamics \cite{CAI2023103226,Lemenkova202405,XIONG2022109132}. However, the question arises regarding the choice of classification algorithms and \hl{efficiency of }software for data processing\hl{. }The integration of \hl{RS} data with \hl{Geographic Information System }(GIS) for environmental mapping has been the subject of research for decades and continues to be a challenging topic. Many approaches to the problem of satellite image analysis have been proposed since the development of satellite missions \cite{NOY2023103377,Lemenkova202403,Su02062020,Potter18102014,Lemenkova202402}. \hl{A new line of research  arises at the intersection between programming methods of Machine Learning (ML) and the spatial analysis using satellite images. The application of such tools and methods enables to process images more effectively and accurately, as reported in previous studies} \cite{rs16244812,w17010050,fire8010008}. 

\hl{Hence}, the performance of image classification methods is linked to the quality of \hl{RS} data. In the field of cartography, Landsat 8-9 \hl{Operational Land Imager and Thermal Infrared Sensor (OLI/TIRS)} image data open new opportunities since the spatial resolution of these data is an advantage for determining land use classes. In this context, the launch of new satellite constellations of Landsat 8-9 OLI/TIRS as \hl{EO} systems allows the acquisition of time series. The acquisition of these data opens new perspectives in cartography \cite{Farah05012010,Lemenkova202322,Layati03042022,Layati28082024}.

%----------- section ----------->
\section{Objective and motivation}

This paper focuses specifically on the application of \hl{ML} methods to satellite image processing incorporated into the GRASS GIS software through the embedded Python Scikit-Learning libraries. \hl{Using ML in the traditional Geographic Information System (GIS) for mapping land cover types is a challenge that has been researched for many years and has many proposed solutions}. Machine learning is an artificial approach to geospatial data processing that relies on mathematical and statistical algorithms to give computers the ability to “learn” from data. \hl{In this context, ML} can be applied to satellite images such as Landsat 8-9 OLI/TIRS since pixels are discrete quantitative variables \cite{DAPAZGOMES,Lemenkova202026,LAROCHEPINEL2024109163}. 

Due to \hl{specific advantages of ML}, the functions of the programs improve their performance in solving geospatial analysis tasks \cite{More02072020,Lemenkova202403,GarciaPedrero03042017}. In the field of satellite image processing, this concerns the usual image processing tasks such as pixel recognition \cite{NASIRI2023103555,EBRAHIMY2023103390}, spectral reflectance analysis \cite{FU2024104421}, classification optimization \cite{PADHY2024100278,Lemenkova202401}, the development of more advanced methods of image analysis and interpretation, as well as the implementation of these methods to cartographic visualization, as discussed earlier \cite{HABIB20231,Lemenkova202406,KHALDI2024104191}. Compared to traditional approaches, the improvements brought by \hl{the ML} algorithms concern the \hl{modelling} of data with an average statistical error as low as possible.

The \hl{Geographic Resources Analysis Support System} (GRASS) GIS software was chosen as the satellite image processing tool for its combination of efficiency and cross-platform compatibility. This paper uses the GRASS \hl{GIS ML} modules ('r.learn.train', 'r.random', 'r.learn.predict' and 'r.category') with the 'ExtraTreesClassifier' and 'DecisionTreeClassifier' models for supervised applications, training and classification of six satellite images. The results of the \hl{ML} approach to image processing are compared with the results of unsupervised classification performed using the GRASS GIS modules 'i.cluster' and 'i.maxlik'. \hl{The scripts of this software are similar to libraries in programming languages.}

In programming libraries, image processing algorithms are implemented as a set of functions used to process raster files, perform geospatial analysis, and create maps in the form of visualized images \cite{Lemenkova202227,PEREIRA2023106653,Lemenkova202229,ESKANDARI2024}. This is the fundamental approach that differentiates the method of using script libraries from traditional GIS. The use of programming approaches to processing and \hl{analysing} satellite images allows handling complex representations of landscape features and detected objects \hl{recognised} as land cover types that have various properties (extent and heterogeneity of landscapes, its texture, size of plots, fragmentation and topology, etc.).

\hl{As well as serving as the first descriptive study of habitat heterogeneity in Rif Mountains, this study also provides guidance for future cartographic mapping using ML modules of GRASS GIS taking into account geographic isolation of diverse landscape patches, heterogeneity and contrast of north African landscapes ranging from arid areas of Sahara to humid coastal zones. }In particular, in this research, an exploration of GRASS GIS \hl{ML} modules is presented and exploited with the aim of demonstrating the applicability of the spatiotemporal \hl{data processing} to environmental \hl{studies in north Africa}.

%----------- section ----------->
\section{Study Area}

The study focuses on northern Morocco, the Rif Mountains, Figure \ref{fig01}. \hl{The Rif} Mountains region is located between latitudes 34\degree 30'-35°15'N and longitudes 4°30'-5°30'W, with an extension of about 30,000 m and an occupation of 4000 km² \cite{ABOUTOFAIL2024112522}. \hl{The} geographical position \hl{of the country }and \hl{its }vast\hl{ }extent\hl{ }result\hl{ }in a great variety of ecosystems. Landscapes vary from humid in the Rif and Atlas, to arid Saharan in the south\hl{, }including sub-humid transition zones and semi-arid in the plains and foothills areas \cite{MCGREGOR20091434}. The Rif Mountains were formed as a result of complex tectonic interactions between the African and Eurasian plates. \hl{ }Complex geological setting of northern Morocco in the Rif region has resulted in moderate tectonic activities \cite{CRINITI2024106861,PARISH199945}. 

\begin{figure}[H]
    \begin{center}
	\hspace{-35pt}\includegraphics[width=14.0cm]{Fig_01}
    \end{center}
    \caption{General topographic map of Morocco with outlined location of the study area. Mapping software: Generic Mapping Tools (GMT) version 6.4.0. Data source: \hl{General Bathymetric Chart of the Oceans} (GEBCO). Map source: author.}
    \label{fig01}
\end{figure}

The geomorphological structures of the region are composed of three main\hl{ }units from north to south: the internal zone, the flysch nappes and the external zone \cite{ELTALIBI2014554,BONARDI2003149}. The crustal thickness beneath the Rif Mountains differs from 36 to 30 km towards the border with the western Mediterranean \cite{MAATE2000409}. \hl{Rif presents a mountain} range that extends across northern Morocco, skirting the Mediterranean Sea and in particular the Alboran Sea on its southern shore. The Rif constitutes a segment of the\hl{ }Alpine ranges. More specifically, it belongs to the Betico-Rifo-Tellian arc in the southwestern Mediterranean \cite{SLIMANI2019101785}. The Rif-Tell belts located in northern Morocco \hl{is} a young Alpine orogenic system\hl{ }that presents\hl{ }metamorphic complexes formed by geodynamic and tectonic activity \cite{MARTINMARTIN2023106367,AITBRAHIM2002187,LEBLANC1984345}. This region is \hl{characterised} by deep sediments\hl{ }related to Tethyan rifting \cite{ARANA1974197,PERRI2022105811,TALMAT2025105515}, which \hl{ }strongly affect soil properties and vegetation types \cite{CHAKKOUR2024,HAMOUCH2024}. \hl{As a result of such complex geologic and geomorphic setting in northern Morocco, }there are very diverse ecosystem types that occupy the landscapes from the foothills, slopes and valleys of limestone mountains to the Saharan sebkhas. The main types of land cover are presented in Figure \ref{fig02}.

\begin{figure}[H]
    \begin{center}
	\hspace{-35pt}\includegraphics[width=14.0cm]{Fig_02}
    \end{center}
    \caption{Land cover types in Morocco. Data source: Food and Agriculture Organization (FAO). The map is created using QGIS software. Map source: author.}
    \label{fig02}
\end{figure}

The geomorphological structure of the Rif can be divided into three segments. The Western Rif is covered with large mountains; its climate is mild, with heavy rains and showers in winter. The Central Rif has diverse reliefs and heterogeneous landscapes. The Eastern Rif is the flattest part of the range \cite{EDDAKIRI2024e02401}. The primary type is presented by the vegetation of the cultivated lands in mosaic distributed along the basement of the Rif Mountains. The vegetation types in the Atlas Mountains region are presented by a closed to open evergreen or semi-deciduous broadleaf forest, with a forest cover of 100 to 40\% for closed to open trees or semi-deciduous trees \cite{BOUBEKRAOUI2023100388}. This geomorphological and climatic heterogeneity corresponds to a bio-ecological and environmental diversity. Thus, the spatial richness of the conditions of climate and geology accentuates this biodiversity and types of land use in Morocco.

The main structure of the Rif is represented by the limestone massifs. The underground water reserves of the Rif maintain the distribution of vegetation and the rivers that originate in\hl{ }Rif \cite{Aqil03042015}. \hl{Due} to these hydrogeological features, short coastal rivers and streams flow from the Rif to the Mediterranean, which are \hl{characterised} by a torrential and erosional regime \cite{Zouhri01062004}. These sources of these rivers supply water to diverse and very heterogeneous types of vegetation around the Rif. These processes briefly illustrate the close link between the geomorphological features, hydrogeology and vegetation of the Rif region \cite{al2017milieux}. 

Seismic,\hl{ }magnetic and thermal activities in the Rif caused by local crustal earthquakes are related to the \hl{geological} instability of this region \cite{DUARTE2024104330}. Specifically, this region presents a southern branch of the Betic-Rif arc that is part of the Alpine ranges. This region has experienced deformations between northwest Africa and southern Spain (Iberia) since the Tertiary \cite{CHABLI2014123}, followed by the convergence between Africa and the Iberian Peninsula \cite{LGHAMOUR2024108941,KAIROUANI2024106762}. \hl{This} tectonic block is \hl{moved} towards the west-southwest and the region still experiences \hl{tectonic} activities \cite{MALUSA2024230533}. At the same time, seismic-tectonic activities shape the regional geomorphology \cite{SADKAOUI2024105219}. In turn, this controls the distribution of land cover types and is expressed in the landscapes. Such links between geology, geomorphology and environmental context have been discussed in previous works \cite{COSTA2024109578,BOUTAHAR2023104828,KAIROUANI2024106762}.

The Mediterranean \hl{coasts} of\hl{ }Morocco is a habitat for numerous \hl{wetland ecosystems including} peatlands with boreal floristic affinities and high conservation value\hl{ } \cite{MULLER202235,AJBILOU2006104}. Occasional patches of rain-fed trees and scrub crops\hl{ }are distributed in the western regions of the country \hl{and depend on water availability} \cite{MARANON1999147}. Rain-fed croplands\hl{, including olive trees and orchards in monoculture systems,} interspersed with areas of sparse vegetation are visible along the Atlantic Ocean coasts \cite{SALMON2015321,ZAYANI2023e22910}. Sparse herbaceous vegetation comprising a complex pattern mixed with shrubs and thickets. Other land cover types are generally presented by urban areas around cities with artificial surfaces and associated areas, and bare soil areas in desert and unoccupied lands, Figure \ref{fig02}. \hl{Besides} natural variability of the Rif \hl{region}, forest areas are affected by human factors \hl{(socio-demographic, economic and urban dynamics)} and hazards \hl{caused by erosion and forest fires} \cite{chebli2018forest}. \hl{The} expansion of agricultural \hl{lands} negatively affect the mountainous forests of the Rif \cite{Adghir02112018,Boutahar01102024,Jimenez01092002}. In addition, the \hl{contrasts from} hot and humid climate,\hl{ varied} topography of the Rif \hl{facilitates} landslides that present geomorphological risks in Northern Morocco \cite{Dahmani10102024,EsSmairi13122022,Manaouch03032024}.

%----------- section ----------->
\section{Materials and Methods}

%----------- subsection ----------->
\subsection{Data}
The Landsat 8-9 OLI/TIRS satellite image data covering the bimonthly period 2023: January, March, May, July, September and November, Figure \ref{fig03}.

\begin{figure}[H]
    \begin{center}
	\hspace{-35pt}\includegraphics[width=12.0cm]{Fig_03}
    \end{center}
    \caption{Landsat images on 2023: (a) – January 26; (b) – March 15; (c) –May 10; (d) – July 29; (e) - September 23; (f) – November 26.}
    \label{fig03}
\end{figure}

The L2 data type is OLI\_TIRS\_L2SP with \hl{the }ground control points \hl{of }version 5. The processing software version is LPGS\_15.3.1c for Landsat-8 scenes and LPGS\_16.3.1 for Landsat-9 (2023) scene. The satellite roll angle is zero for Landsat 8 scenes and -0.001 for Landsat 9. The images were taken during the day at Nadir, with cloud cover less than 10\%. \hl{Here}, only multispectral bands with an acquisition resolution of 30 meters were used, except for panchromatic bands with a resolution of 15 m. The combination of multispectral bands is mainly used for the creation of false/true color composites for classification, especially for image visualization. The available data were listed using the search template in "g.list rast". 

%----------- subsection ----------->
\subsection{Workflow}

\hl{We first plot the workflow obtained from each methodological step to provide a general picture of the approach used in this study. Since various GIS include different steps of data processing, so we next summarise the scheme into 12 well-recognised research steps, based on modules of GRASS GIS for geospatial data processing, Figure} \ref{fig04}. 

\begin{figure}[H]
    \begin{center}
	\hspace{-35pt}\includegraphics[width=13.0cm]{Fig_04}
    \end{center}
    \caption{\hl{Workflow diagram illustrating major steps of research design used for organisation of data processing tasks. Software: GIMP, version 2.10. Diagram source: author.}}
    \label{fig04}
\end{figure}

\hl{Major research steps include the following points: 1) Data collection; 2) Software selection; 3) Quality control; 4) Data pre-processing; 5) Image processing 6) Clustering (unsupervised classification); 7) Machine learning; 8) Training dataset; 9) Supervised classification; 10) Numerical computation; 11) Cartographic mapping; 12) Validation.}

\hl{The data preprocessing procedures included several steps which include geometric and atmospheric correction for optical remote sensing images. The geometric correction was performed using 'i.rectify' module of GRASS GIS which corrected each image by performing a coordinate transformation for the pixels in the images using control points. Hence the x,y cell coordinates on the images were re-calculated to a transformation matrix and then converted to standard map coordinates for each pixel in the image. Following that, the atmospheric correction was performed using the 6S algorithm embedded in the 'i.atcorr' module of GRASS GIS. The atmospheric correction is essential to evaluating numerically the amount of wavelengths reflected of diverse land cover types of the surface in various seasonal periods} \cite{f14071414,ABDELRAHMANHEGAB2024716}. \hl{The pansharpening has been performed using the 'i.pansharpen' module of GRSAS GIS which uses image fusion algorithms to sharpen multispectral channels with high-resolution panchromatic channels. In this way, the resolution of the multispectral is improved. The effectiveness of the pansharpening is described in earlier studies on RS data processing} \cite{rs15082017,rs13173433}.

Mapping of topographic data was carried out using GMT using existing methods \cite{Lemenkova202221,BECKER20051075,Lemenkova202217,10537274,5369602}, while mapping of land use data was carried out using QGIS \cite{FANG202031}. Mapping of land cover maps from \hl{RS} data is based on satellite image classification methods.

The aim of this approach is to model the data and predict for each pixel a label, i.e. Land Cover Class. Using \hl{ML} methods, these methods are divided into two categories: supervised classification and unsupervised classification. This study involves the use and comparison of these two methods. The image processing process is composed of the following steps, implemented in the form of free and open source modules on GRASS \hl{GIS}. The code for organising the project and importing data is presented below.

\begin{footnotesize}
\begin{verbatim}
# Create the project with new location from raster map 
(file must contain projection metadata):
# grass -c myraster.tif /home/user/grassdata/mynewlocation
grass
#cd /Users/polinalemenkova/grassdata
#grass -c LC09_L2SP_179073_20220419_20230421_02_T1_SR_B1.tif 
/Users/polinalemenkova/grassdata/Morocco
# ----IMPORT AND PREPROCESSING-------------------------->
# g.mapset location=Morocco mapset=PERMANENT
g.list rast
# importing the image subset with 7 Landsat bands and display the raster map
r.import input=/Users/polinalemenkova/grassdata/Morocco/
LC09_L2SP_201036_20230126_20230313_02_T1_SR_B1.TIF 
output=L9_2023_J_01 extent=region resolution=region --overwrite
r.import input=/Users/polinalemenkova/grassdata/Morocco/
LC09_L2SP_201036_20230126_20230313_02_T1_SR_B2.TIF 
output=L9_2023_J_02 extent=region resolution=region
r.import input=/Users/polinalemenkova/grassdata/Morocco/
LC09_L2SP_201036_20230126_20230313_02_T1_SR_B3.TIF 
output=L9_2023_J_03 extent=region resolution=region
r.import input=/Users/polinalemenkova/grassdata/Morocco/
LC09_L2SP_201036_20230126_20230313_02_T1_SR_B4.TIF
 output=L9_2023_J_04 extent=region resolution=region
r.import input=/Users/polinalemenkova/grassdata/Morocco/
LC09_L2SP_201036_20230126_20230313_02_T1_SR_B5.TIF
 output=L9_2023_J_05 extent=region resolution=region
r.import input=/Users/polinalemenkova/grassdata/Morocco/
LC09_L2SP_201036_20230126_20230313_02_T1_SR_B6.TIF 
output=L9_2023_J_06 extent=region resolution=region
r.import input=/Users/polinalemenkova/grassdata/Morocco/
LC09_L2SP_201036_20230126_20230313_02_T1_SR_B7.TIF
 output=L9_2023_J_07 extent=region resolution=region
#
g.list rast
# g.remove -f type=raster name=L9_2023_J_J_01
\end{verbatim}
\end{footnotesize}

First, the import and preprocessing of six Landsat images was \hl{implemented} using the modules 'r.import' that reads the data\hl{. Then, the }'i.landsat.toar' \hl{was used to convert} Landsat 8-9 OLI/TIRS images and \hl{to transform} the calibrated digital numbers for the reflectance at the top of the atmosphere. In this way, the quality of the images \hl{has been} improved. On the \hl{coloured} composition in "false colors", the vegetation appears in different shades depending on the species, but also on the environmental conditions: geomorphology of the mountain slope, the curvature and the inclination of the slope that affect the shadows on the images, Figure \ref{fig05}.

\begin{figure}[H]
    \begin{center}
	\hspace{-35pt}\includegraphics[width=14.0cm]{Fig_05}
    \end{center}
    \caption{Landsat 8-9 multispectral images: false and true colour composites.}
    \label{fig05}
\end{figure}

The spectral combination calculation integrates different spectral bands of Landsat based on the physical properties of the land covers reflected in the images. Therefore, the spectral bands are useful to improve the discrimination between landscape classes, Figure \ref{fig05}. The vegetation observed on the Atlantic Ocean side appears in darker reds than the surrounding vegetation around the Rif Mountains. The water that absorbs almost all wavelengths is very dark, up to black, while the empty ground surface appears very light, in tones ranging from beige to ochre, Figure \ref{fig05}.

Creation of \hl{colour} composites (Figure \ref{fig05}) was performed based on the Landsat multispectral bands using the 'r.composite' module and visualization of the images using the 'd.rast' module. The estimation of the color composition model aims to better \hl{recognise }the ten land cover classes representing landscape patches in the Rif Mountains and its surroundings in the images. \hl{Modelling} and simulation of the images using the principle of \hl{colour} compositions includes 'true color' and 'false color' images, Figure \ref{fig05}. Here, the spectral bands of the images are visualized in several \hl{colour} combinations that are widely used in \hl{RS} because they are very suitable for studying vegetation and landscape. For example, the combination with band number 5 (infrared, IR) is based on the properties of vegetation that reflect very strongly near-IR radiation.

\begin{footnotesize}
\begin{verbatim}
# ----CREATING COLOR COMPOSITES----->
# false color
r.composite blue=L9_2023_J_07 green=L9_2023_J_05 red=L9_2023_J_03
 output=L9_2023_J_753 --overwrite
d.mon wx0
d.rast L9_2023_J_753
d.out.file output=Morocco_753 format=jpg --overwrite
# false color: NIR band B05 in the red channel, red band B04 in the green 
channel and green band B03 in the blue channel
r.composite blue=L9_2023_J_03 green=L9_2023_J_04 red=L9_2023_J_05
 output=L9_2023_J_345 --overwrite
d.mon wx0
d.rast L9_2023_J_345
d.out.file output=Morocco_345 format=jpg --overwrite
# true color
r.composite blue=L9_2023_J_02 green=L9_2023_J_03 red=L9_2023_J_04
output=L9_2023_J_234 --overwrite
d.mon wx0
d.rast L9_2023_J_234
d.out.file output=Morocco_J_234 format=jpg --overwrite
\end{verbatim}
\end{footnotesize}

\hl{The use of ML approach to} satellite image processing by GRASS GIS consists of two phases. The first consists of estimating a model from data by automatic clustering. The unsupervised clustering approach groups comparable samples within the same class using the k-means algorithm. The image clustering and classification according to the maximum likelihood algorithm (unsupervised learning) using the 'i.cluster' and 'i.maxlik' modules. The groups, i.e. the clusters are made up of pixels with similar spectral values. The unsupervised classification uses the clustering approach performed using the 'i.cluster' module as explained above.

\hl{The }unsupervised classification \hl{employs} clustering approach performed using the "i.cluster" module. The "i.cluster" generates a signature file and reports the cluster maps, or regions, that were automatically generated using the "k-means" data partitioning algorithm. The GRASS SIG code used for this is the following (example for the 2013 image, repeated for images from all other years using the same technique): "i.cluster group=L8\_2013 subgroup=res\_30m signaturefile=cluster\_L8\_2013 classes=8 reportfile=rep\_clust\_L8\_2013.txt". Therefore, samples within a class are dissimilar from samples in others.

\begin{footnotesize}
\begin{verbatim}
# ---Clustering and Classification ---->
# grouping data by i.group
# Set computational region to match the scene
g.region raster=L9_2023_J_01 -p
i.group group=L9_2023_J subgroup=res_30m \
input=L9_2023_J_01,L9_2023_J_02,L9_2023_J_03,L9_2023_J_04,
L9_2023_J_05,L9_2023_J_06,L9_2023_J_07 --overwrite
# Clustering: generating signature file and report using 
k-means clustering algorithm
i.cluster group=L9_2023_J subgroup=res_30m \
  signaturefile=cluster_L9_2023_J \
  classes=10 reportfile=rep_clust_L9_2023_J.txt --overwrite
# Classification by i.maxlik module
i.maxlik group=L9_2023_J subgroup=res_30m \
  signaturefile=cluster_L9_2023_J \
  output=L9_2023_J_cluster_classes reject=L9_2023_J_cluster_reject 
# Mapping
d.mon wx0
g.region raster=L9_2023_J_cluster_classes -p
r.colors L9_2023_J_cluster_classes color=bcyr
d.rast L9_2023_J_cluster_classes
d.legend raster=L9_2023_J_cluster_classes title="26 January 2023" 
title_fontsize=14 font="Helvetica" fontsize=12 bgcolor=white
 border_color=white
d.out.file output=Morocco_2023_Jan format=jpg --overwrite
# Mapping rejection probability
d.mon wx1
g.region raster=L9_2023_J_cluster_classes -p
r.colors L9_2023_J_cluster_reject color=wave -e
d.rast L9_2023_J_cluster_reject
d.legend raster=L9_2023_J_cluster_reject title="26 January 2023" 
title_fontsize=14 font="Helvetica" fontsize=12 bgcolor=white 
border_color=white
d.out.file output=Morocco_2023_reject format=jpg --overwrite
\end{verbatim}
\end{footnotesize}

In this way, each land cover class is associated with each cluster. Calculating the pixel rejection probability using accuracy assessment for quality control. With the unsupervised modeling done, the new data obtained is submitted as the training dataset for \hl{ML} to obtain better image classification results. The clustering results are used as observation data available for supervised learning. Thus, the results of this step are used for \hl{ML} as the training dataset.

Subsequently, \hl{ML} processing of images was performed using supervised classification that uses a sequence of modules. The ability of \hl{ML} to model geospatial data relies heavily on the acquisition of training data as the algorithm feeds on the input dataset, so this is an important step. \hl{ML} for image classification using “r.learn.train” and “r.learn.predict” models.

The algorithm is based on the use of the ‘decision tree’ and ‘extra trees’ classifiers that are integrated into the GRASS GIS module ‘r.learn.train’. Here, the data are labeled, i.e. This is a supervised learning with the technical specificities mentioned in the code below. First, the ‘r.random’ module was used to generate training pixels from a previous classification of the land cover of Morocco. All these algorithms were implemented using the GRASS GIS module ‘r.learn.train’ using the programming code presented below.

\begin{footnotesize}
\begin{verbatim}
# MACHINE LEARNING ------>
# g.list rast
g.region raster=L9_2023_J_01 -p
# First, we are going to generate some training pixels from an older (1996) 
land cover classification:
r.random input=L9_2023_cluster_classes seed=100 npoints=1000 
raster=training_pixels --overwrite
# Next, we create the imagery group with all 
Landsat-8 OLI/TIRS 7 (2000) bands:
i.group group=L9_2023_J 
input=L9_2023_J_01,L9_2023_J_02,L9_2023_J_03,L9_2023_J_04,
L9_2023_J_05,L9_2023_J_06,L9_2023_J_07 --overwrite
# Then use these training pixels to perform a classification on 
recent Landsat - 2022 image:
#
# train a decision tree classification model using r.learn.train
r.learn.train group=L9_2023_J training_map=training_pixels \
    model_name=DecisionTreeClassifier n_estimators=500 
    save_model=rf_model.gz --overwrite
# perform prediction using r.learn.predict
r.learn.predict group=L9_2023_J load_model=rf_model.gz
 output=rf_classification --overwrite
# check raster categories - they are automatically applied 
to the classification output
r.category rf_classification
# copy color scheme from landclass training map to result
r.colors rf_classification raster=training_pixels
# display
d.mon wx0
d.rast rf_classification
r.colors rf_classification color=roygbiv -e
d.legend raster=rf_classification title="Decision Tree: 01/2023" 
title_fontsize=14 font="Helvetica" fontsize=12 bgcolor=white 
border_color=white
d.out.file output=DT_2023_01 format=jpg --overwrite
\end{verbatim}
\end{footnotesize}

In the next step, the prediction was performed using the module "r.learn.predict" which was performed in the code: "r.learn.predict group=L8\_2013 load\_model=rf\_model.gz output=rf\_classification". The raster categories were examined using the module "r.category" of GRASS SIG, based on the data automatically applied to the classification output. For this, the following code was used: "r.category rf\_classification". The maps were visualized using the combination of two main modules: "d.rast" which visualizes the map itself and "d.legend" which adds the legend on this map. The images were converted to bitmap format using the module "d.out.file", e.g., "d.out.file output=Maroc\_2013 format=jpg".

\begin{footnotesize}
\begin{verbatim}
# train a extra trees classification model using r.learn.train
r.learn.train group=L9_2023_J training_map=training_pixels \
   model_name=ExtraTreesClassifier n_estimators=500 
   save_model=rf_model.gz --overwrite
# perform prediction using r.learn.predict
r.learn.predict group=L9_2023_J load_model=rf_model.gz  output=rf_classification --overwrite
# check raster categories applied to the classification output
r.category rf_classification
# copy color scheme from landclass training map to result
r.colors rf_classification raster=training_pixels
# display
d.mon wx0
d.rast rf_classification
r.colors rf_classification color=bgyr -e
d.legend raster=rf_classification title="Extra Tree: 01/2023" 
title_fontsize=14 font="Helvetica" fontsize=12 bgcolor=white 
border_color=white
d.out.file output=ET_2023_01 format=jpg --overwrite
\end{verbatim}
\end{footnotesize}

This was used as a basis for \hl{ML} using the following code: “r.random input=L8\_2013\_cluster\_classes seed=100 npoints=1000 raster=training\_pixels”. After that, the “r.learn.train” module was used to train a module using the training map that was created in the previous step by applying the built-in algorithm. The code used is as follows: “r.learn.train group=L8\_2013 training\_map=t\_pixels model\_name=DecisionTreeClassifier n\_estimators=500 save\_model=rf\_model.gz”. For all approaches, the model name was changed using the “model\_name” function and the name of the model in question, e.g., “modelage=RandomForestRegressor”.

Since the pixel labels are discrete (not continuous which requires the use of regression), the classification methods 'decision trees' and 'supplementary trees' were chosen. These \hl{ML} algorithms predict the probability of pixel values suitable for the target class on the image. Therefore, the machine performs a classification of the image pixels into 10 land cover classes on the Rif Mountains. The categorization of the classes was done using the 'r.category' module of GRASS GIS. 

Finally, the visualization of images using the auxiliary graphic modules of GRASS GIS for cartographic processing: 'd.mon', 'd.rast', 'r.colors' and 'd.legend'. The images were saved as bitmap files using the module 'd.out.file'. The calculation and estimation of classification accuracy are done by the function 'reject' for the evaluation of the probability of the percentage of correctly classified pixels on the image. The evaluation of image processing quality and algorithm performance are done by GRASS GIS.

%----------- section ----------->
\section{Results and Discussion}

\hl{Mountain areas in northern Morocco are particularly sensitive to climate change and pastoral activities. Seasonal variations in climate and effects from human activities strongly influence agro-ecosystems and cause changes in land cover types. Although there are some efforts on mapping the landscape dynamics in such areas, any seasonal response of the land cover types will is better evaluated using quantitative assessment of RS data and ML methods. Using examples from the Rif Atlas in Morocco, this paper explored the seasonal variations in land cover types in northern Morocco in 2023 caused by the effects from livestock management, arboriculture and climate change. The selected period was assessed for the year 2023 using the available time series from Landsat dataset collection. The results} of satellite image classification \hl{were obtained} using three approaches: unsupervised classification and supervised classification by \hl{ML} methods \hl{which employs} two \hl{distinct algorithms:} "Decision tree classifier" and "Extra tree classifier". Figure \ref{fig06} shows the images generated using the clustering classification method, which illustrate the land cover changes with a bi-monthly time interval, Table \ref{tab01}.

\begin{table}[H]
\centering
\begin{tabularx}{\textwidth}{| p{1.5cm} | X | X | X | X | X | X | X | X | X | X | }
  \toprule
   Class & 1 & 2 & 3 & 4 & 5 & 6 & 7 & 8 & 9 & 10 \\
   \midrule
    January & 766 & 465 & 554 & 736 & 1031 & 785 & 462 & 1016 & 646 & 149 \\
    March & 720 & 341 & 696 & 673 & 1060 & 791 & 587 & 821 & 720 & 203 \\
    May & 742 & 348 & 545 & 449 & 904 & 912 & 905 & 722 & 645 & 501 \\
    July & 766 & 465 & 554 & 736 & 1031 & 785 & 462 & 1016 & 646 & 149 \\
    September & 668 & 424 & 674 & 318 & 793 & 791 & 1099 & 1023 & 683 & 222 \\
    November & 781 & 483 & 701 & 675 & 1097 & 429 & 744 & 833 & 734 & 225 \\
    \bottomrule
\end{tabularx}
\caption{Bimonthly variations of in Rif Mountains by pixels over 10 land cover classes in 2023 in two (cold/hot) different seasons: 1) January, March and May; 2) July, September and November}\label{tab01}
\end{table}

Figure \ref{fig07} shows that the classification using the “decision tree classifier” method improved the results since the \hl{ML} results are more robust and stable compared to the unsupervised classification. The January scene shows the least plant activity which gradually increases from March to May and decreases respectively after the hot summer period until November.

\begin{figure}[H]
    \begin{center}
	\hspace{-35pt}\includegraphics[width=14.0cm]{Fig_06}
    \end{center}
    \caption{Satellite image classification using unsupervised classification clustering methods and the GRASS GIS 'MaxLike' approach.}
    \label{fig06}
\end{figure}

On the one hand, vegetation varied seasonally in the Rif region according to the natural dynamics of the phenological cycle during different calendar months. However, compared to other land cover types, the mosaic vegetation types of deciduous forests developed partially and differently for both algorithms, with more density and a homogeneous pattern across the entire study area.

\begin{figure}[H]
    \begin{center}
	\hspace{-35pt}\includegraphics[width=14.0cm]{Fig_07}
    \end{center}
    \caption{Results of satellite image classification using supervised \hl{ML} method with GRASS GIS 'Decision Tree' algorithm.}
    \label{fig07}
\end{figure}

On the other hand, technical improvements of the parameters that differ in various algorithms ‘MaxLike’, ‘DecisionTreeClassifier’ and ‘ExtraTreesClassifier’ allow to better visualize land cover classes on the images. Thus, the most complex landscape patches with high heterogeneity can be better distinguished, including sparsely distributed and densely distributed vegetation segments. Here, the rejection threshold map contains the index of a confidence level calculated for each classified pixel in the classified image of the 10 land cover classes in the Morocco region. Confidence intervals are defined and the rejection map is visualized with lower values meaning 'keep the correctly classified pixel' and higher values meaning “reject a misclassified pixel that could belong to another class category”. Thus, this map identifies cells in the classified raster image that have a low probability, i.e. a high rejection index, of being correctly assigned to the target class.

\begin{equ}[!ht]
\begin{equation}
\scriptsize
\begin{vmatrix}
   & 1 & 2 & 3 & 4 & 5 & 6 & 7 & 8 & 9 & 10 \\
1 & 0 &  &  &  &  &  &  &  &  &  \\
2 & 1.7 & 0 &  &  &  &  &  &  &  &  \\
3 & 3.3 & 1.1 & 0 &  &  &  &  &  &  &  \\
4 & 3.3 & 1.0 & 0.7 & 0 &  &  &  &  &  &  \\
5 & 4.2 & 1.9 & 0.9 & 1.2 & 0 &  &  &  &  &  \\
6 & 3.9 & 1.6 & 1.2 & 0.8 & 1.0 & 0 &  &  &  &  \\  
7 & 3.5 & 1.8 & 1.6 & 1.2 & 1.2 & 0.6 & 0 &  &  &  \\  
8 & 4.6 & 2.3 & 1.5 & 1.7 & 0.7 & 1.2 & 1.1 & 0 &  &  \\
9 & 5.0 & 3.0 & 2.2 & 2.5 & 1.5 & 1.9 & 1.6 & 0.9 & 0 &  \\   
10 & 3.5 & 2.6 & 2.2 & 2.3 & 1.8 & 2.0 & 1.7 & 1.4 & 0.9 & 0 \\          
\end{vmatrix}
 %
\hspace{3pt}
\begin{vmatrix}
    & 1 & 2 & 3 & 4 & 5 & 6 & 7 & 8 & 9 & 10 \\
1 & 0 &  &  &  &  &  &  &  &  &  \\
2 & 2.5 & 0 &  &  &  &  &  &  &  &  \\
3 & 3.0 & 1.1 & 0 &  &  &  &  &  &  &  \\
4 & 3.5 & 1.3 & 0.6 & 0 &  &  &  &  &  &  \\
5 & 4.3 & 2.2 & 0.9 & 0.9 & 0 &  &  &  &  &  \\
6 & 4.7 & 2.6 & 1.3 & 1.1 & 0.7 & 0 &  &  &  &  \\  
7 & 5.1 & 3.1 & 1.7 & 1.7 & 1.0 & 0.8 & 0 &  &  &  \\  
8 & 5.4 & 3.8 & 2.3 & 2.3 & 1.7 & 1.1 & 1.0 & 0 &  &  \\
9 & 5.7 & 3.8 & 2.5 & 2.4 & 1.9 & 1.5 & 1.0 & 0.8 & 0 &  \\   
10 & 5.5 & 4.2 & 3.0 & 2.9 & 2.6 & 2.0 & 1.9 & 1.1 & 1.1 & 0 \\              
\end{vmatrix} 
\end{equation}
\caption{Bimonthly class separability matrices between 10 classes of land cover types in Rif Mountains, Morocco: January and March 2023 (1)}
\end{equ}

\begin{equ}[!ht]
\begin{equation}
\scriptsize
\begin{vmatrix}
   & 1 & 2 & 3 & 4 & 5 & 6 & 7 & 8 & 9 & 10 \\
1 & 0 &  &  &  &  &  &  &  &  &  \\
2 & 2.6 & 0 &  &  &  &  &  &  &  &  \\
3 & 4.0 & 0.9 & 0 &  &  &  &  &  &  &  \\
4 & 4.5 & 1.7 & 0.9 & 0 &  &  &  &  &  &  \\
5 & 4.4 & 1.3 & 0.7 & 1.1 & 0 &  &  &  &  &  \\
6 & 4.0 & 1.5 & 1.2 & 1.7 & 0.8 & 0 &  &  &  &  \\  
7 & 5.0 & 2.1 & 1.4 & 1.0 & 1.0 & 1.2 & 0 &  &  &  \\  
8 & 5.5 & 2.2 & 1.4 & 0.6 & 1.3 & 1.9 & 0.7 & 0 &  &  \\
9 & 5.7 & 3.0 & 2.2 & 1.4 & 2.2 & 2.6 & 1.2 & 0.9 & 0 &  \\   
10 & 5.2 & 3.3 & 2.7 & 2.2 & 2.6 & 2.9 & 1.9 & 1.7 & 1.1 & 0 \\          
\end{vmatrix}
 %
\hspace{3pt}
\begin{vmatrix}
    & 1 & 2 & 3 & 4 & 5 & 6 & 7 & 8 & 9 & 10 \\
1 & 0 &  &  &  &  &  &  &  &  &  \\
2 & 1.7 & 0 &  &  &  &  &  &  &  &  \\
3 & 3.3 & 1.1 & 0 &  &  &  &  &  &  &  \\
4 & 3.3 & 1.0 & 0.7 & 0 &  &  &  &  &  &  \\
5 & 4.2 & 1.9 & 0.9 & 1.2 & 0 &  &  &  &  &  \\
6 & 3.9 & 1.6 & 1.2 & 0.8 & 1.0 & 0 &  &  &  &  \\  
7 & 3.5 & 1.8 & 1.6 & 1.2 & 1.2 & 0.6 & 0 &  &  &  \\  
8 & 4.6 & 2.3 & 1.5 & 1.7 & 0.7 & 1.2 & 1.1 & 0 &  &  \\
9 & 5.0 & 3.0 & 2.2 & 2.5 & 1.5 & 1.9 & 1.6 & 0.9 & 0 &  \\   
10 & 3.5 & 2.6 & 2.2 & 2.3 & 1.8 & 2.0 & 1.7 & 1.4 & 0.9 & 0 \\              
\end{vmatrix} 
\end{equation}
\caption{Bimonthly class separability matrices between 10 classes of land cover types in Rif Mountains, Morocco: May and July 2023 (2).}
\end{equ}

\begin{equ}[!ht]
\begin{equation}
\scriptsize
\begin{vmatrix}
   & 1 & 2 & 3 & 4 & 5 & 6 & 7 & 8 & 9 & 10 \\
1 & 0 &  &  &  &  &  &  &  &  &  \\
2 & 2.8 & 0 &  &  &  &  &  &  &  &  \\
3 & 4.8 & 1.1 & 0 &  &  &  &  &  &  &  \\
4 & 3.6 & 1.0 & 0.6 & 0 &  &  &  &  &  &  \\
5 & 6.3 & 2.1 & 1.0 & 1.1 & 0 &  &  &  &  &  \\
6 & 6.8 & 2.3 & 1.3 & 1.1 & 0.6 & 0 &  &  &  &  \\  
7 & 8.2 & 3.1 & 2.1 & 1.7 & 1.1 & 0.8 & 0 &  &  &  \\  
8 & 8.6 & 3.6 & 2.6 & 2.1 & 1.8 & 1.4 & 0.7 & 0 &  &  \\
9 & 8.6 & 4.0 & 3.2 & 2.5 & 2.5 & 2.1 & 1.5 & 0.8 & 0 &  \\   
10 & 7.1 & 4.1 & 3.4 & 2.9 & 2.9 & 2.6 & 2.2 & 1.6 & 1.0 & 0 \\          
\end{vmatrix}
 %
\hspace{3pt}
\begin{vmatrix}
    & 1 & 2 & 3 & 4 & 5 & 6 & 7 & 8 & 9 & 10 \\
1 & 0 &  &  &  &  &  &  &  &  &  \\
2 & 2.0 & 0 &  &  &  &  &  &  &  &  \\
3 & 3.3 & 1.0 & 0 &  &  &  &  &  &  &  \\
4 & 3.9 & 1.6 & 0.9 & 0 &  &  &  &  &  &  \\
5 & 5.0 & 2.3 & 1.3 & 0.8 & 0 &  &  &  &  &  \\
6 & 3.6 & 1.6 & 0.8 & 1.0 & 0.9 & 0 &  &  &  &  \\  
7 & 5.5 & 2.8 & 1.8 & 1.6 & 0.8 & 1.0 & 0 &  &  &  \\  
8 & 6.0 & 3.2 & 2.2 & 1.8 & 1.0 & 1.5 & 0.7 & 0 &  &  \\
9 & 6.1 & 3.6 & 2.8 & 2.6 & 1.8 & 1.9 & 1.3 & 1.0 & 0 &  \\   
10 & 5.3 & 3.8 & 3.2 & 3.1 & 2.6 & 2.5 & 2.2 & 2.0 & 1.2 & 0 \\              
\end{vmatrix} 
\end{equation}
\caption{Bimonthly class separability matrices between 10 classes of land cover types in Rif Mountains, Morocco: September and November 2023 (3).}
\end{equ}

Figure 7 shows the land cover type patterns over the Rif Mountains and surrounding areas, defined using an improved \hl{ML} algorithm.

\begin{figure}[H]
    \begin{center}
	\hspace{-35pt}\includegraphics[width=14.0cm]{Fig_08}
    \end{center}
    \caption{Results of satellite image classification using supervised \hl{ML} method with GRASS GIS 'Extra Tree' algorithm.}
    \label{fig08}
\end{figure}

Compared to the previous results in Figure \ref{fig06} and Figure \ref{fig07}, the results of the Extra tree are not quite similar for the classification maps based on the “Decision tree” (Figure \ref{fig07}) and “Extra tree” algorithms (Figure \ref{fig08}), because they rely on the extremely random trees method. The “Extra tree” algorithm is similar to the random forest algorithm in its nature, but it processes images much faster. Therefore, technically, this algorithm is an improved ensemble supervised \hl{ML} method.

However, this algorithm creates many decision trees using random sampling of pixels for each tree. The most important and unique feature of the “Extra trees” algorithm is that it randomly selects a pixel division value. Therefore, compared with the “decision tree” algorithms, the land cover types are diverse and uncorrelated. The best performance among the three image processing methods is demonstrated by the “decision tree” algorithms which divide the image into land cover classes with precision and detail. The accuracy assessment is presented in Figure \ref{fig09} and computational results available in Table \ref{tab03} and Table \ref{tab04}.

\begin{longtable}{@{\extracolsep{\fill}}p{1.5cm}cccccccc@{}}
\caption{Bimonthly evaluation of the class means of Digital Numbers (DN) computed for pixels assigned to 10 land cover classes in landscapes of Rif Mountains, north Morocco. Evaluated each multispectral band (1 to 7) of Landsat 8-9 OLI/TIRS. Period: 2023, January to May.} \label{tab03}\\
\hline \textbf{Class} & \textbf{Band 1} & \textbf{Band 2} & \textbf{Band 3} & \textbf{Band 4} & \textbf{Band 5} & \textbf{Band 6} & \textbf{Band 7} \\ \hline 
\endfirsthead

\multicolumn{5}{c}%
{{\tablename\ \thetable{} -- continued from previous page}} \\
\textbf{Class} & \textbf{Band 1} & \textbf{Band 2} & \textbf{Band 3} & \textbf{Band 4} & \textbf{Band 5} & \textbf{Band 6} & \textbf{Band 7}  \\ \hline 
\endhead

\hline \multicolumn{8}{r}{{Continued on next page}} \\ \hline
\endfoot

\hline \hline
\endlastfoot

%\bf{2019} &  &  &  &  &  &  & \\
\multicolumn{8}{| c |}{\textbf{2023: January}}\\
\hline
1 & 8053.66 & 8105.74 & 7961.36 & 7451.51 & 7386.15 & 7411.52 & 7385.87 \\
2 & 7529.44 & 7754.46 & 8346.18 & 8318.72 & 11431.1 & 10065.3 & 8996.52 \\
3 & 8054.51 & 8324.72 & 9136.1 & 9317.05 & 12308.5 & 12725.3 & 11052.1 \\
4 & 7889.8 & 8102.02 & 8942.29 & 8745.57 & 14901.1 & 12282.7 & 10135 \\
5 & 8332 & 8667.33 & 9728.08 & 10061.3 & 14289.8 & 14818.8 & 12637.9 \\
6 & 8110.13 & 8324.27 & 9467.21 & 9012.78 & X17982.6& 13440.9 & 10707.4 \\
7 & 8251.2 & 8447.13 & 9831.94 & 9044.77 & 21802.8 & 13717.8 & 10608.9 \\
8 & 8640.53 & 9039.94 & 10407.5 & 10794.4 & 16840.5 & 16286.3 & 13653.7 \\
9 & 9111.71 & 9625.79 & 11316.9 & 12294.1 & 17724.7 & 19083.2 & 16430 \\
10 & 11270 & 12124 & 14516 & 16218.7 & 21505.9 & 23336.4 & 20483 \\
\hline
%\bf{2022} &  &  &  & & & & \\
\multicolumn{8}{| c |}{\textbf{2023: March}}\\
\hline
1 & 7732.78 & 7928.37 & 7997.69 & 7540.35 & 7530.62 & 7850.57 & 7839.92 \\
2 & 7850.68 & 8119.44 & 9096.92 & 8986.22 & 16313.3 & 12929.6 & 10370.1 \\
3 & 8406.34 & 8821.78 & 10011.9 & 10547.6 & 15377.5 & 14966.2 & 12509.7 \\
4 & 8377.33 & 8825.2 & 10246.8 & 10459 & 19450.9 & 15628.6 & 12195.9 \\
5 & 8777.26 & 9318.35 & 10778.5 & 11790.6 & 16488.5 & 17202.7 & 14298.8 \\
6 & 9050.82 & 9764.11 & 11495.9 & 12894.6 & 18940.7 & 17954 & 14120.4 \\
7 & 9129.74 & 9761.09 & 11407.2 & 12937.6 & 17616.1 & 19787.9 & 16566.3 \\
8 & 9848.06 & 10823.1 & 12951.4 & 15075.9 & 19854.8 & 20511.1 & 16312.7 \\
9 & 9570.95 & 10302.4 & 12240.4 & 14420.2 & 19123.6 & 22705 & 19360.4 \\
10 & 10759 & 11957.4 & 14571.2 & 17344.9 & 21903.1 & 24000.4 & 19583.5 \\
\hline

%\bf{2023} &  &  &  & & & & \\
\multicolumn{8}{| c |}{\textbf{2023: May}}\\
\hline
1 & 8111.51 & 8199.08 & 8100.88 & 7465.42 & 7287.35 & 7335.55 & 7330.84 \\
2 & 7565.9 & 7794.02 & 8480.23 & 8399.37 & 13105.3 & 11135.9 & 9540.2 \\
3 & 7941.53 & 8209.27 & 9135.12 & 9131.33 & 14663.5 & 13138.1 & 10966.2 \\
4 & 8278.6 & 8635.24 & 9757.47 & 10183.6 & 14710.6 & 15182 & 12997.8 \\
5 & 8007.61 & 8238.4 & 9379.12 & 9007.12 & 18137.6 & 13349.7 & 10649.6 \\
6 & 7959.3 & 8113.15 & 9263.84 & 8488.51 & 22413.9 & 12764.6 & 9831.97 \\
7 & 8460.35 & 8799.31 & 10364.7 & 10092.9 & 20827.9 & 15638.7 & 12351.6 \\
8 & 8552.87 & 8939.44 & 10306.6 & 10660.9 & 17349.2 & 16382.3 & 13624.4 \\
9 & 8999.75 & 9517.87 & 11173.8 & 12159.9 & 17889.8 & 18838.8 & 16259.8 \\
10 & 10309.3 & 11076.4 & 13413.5 & 15174.4 & 20565.5 & 22691.1 & 20127.9 \\
\hline
\end{longtable}

\begin{longtable}{@{\extracolsep{\fill}}p{1.5cm}cccccccc@{}}
\caption{Bimonthly evaluation of the class means of Digital Numbers (DN) computed for pixels assigned to 10 land cover classes in landscapes of Rif Mountains, north Morocco. Evaluated each multispectral band (1 to 7) of Landsat 8-9 OLI/TIRS. Period: 2023, July to November.} \label{tab04}\\
\hline \textbf{Class} & \textbf{Band 1} & \textbf{Band 2} & \textbf{Band 3} & \textbf{Band 4} & \textbf{Band 5} & \textbf{Band 6} & \textbf{Band 7} \\ \hline 
\endfirsthead

\multicolumn{5}{c}%
{{\tablename\ \thetable{} -- continued from previous page}} \\
\textbf{Class} & \textbf{Band 1} & \textbf{Band 2} & \textbf{Band 3} & \textbf{Band 4} & \textbf{Band 5} & \textbf{Band 6} & \textbf{Band 7}  \\ \hline 
\endhead

\hline \multicolumn{8}{r}{{Continued on next page}} \\ \hline
\endfoot

\hline \hline
\endlastfoot

%\bf{2019} &  &  &  &  &  &  & \\
\multicolumn{8}{| c |}{\textbf{2023: July}}\\
\hline
1 & 7500.81 & 7646.92 & 7722.75 & 7376.19 & 7464.38 & 7798.77 & 7755.02 \\
2 & 8127.96 & 8540.32 & 9609.1 & 9705.24 & 15267.4 & 13436.5 & 11015.4 \\
3 & 8252.89 & 8754.25 & 10217.3 & 10220.6 & 19793.6 & 15491.7 & 12155.2 \\
4 & 8632.51 & 9160.58 & 10487.1 & 11145.4 & 15532 & 15921.1 & 13380.2 \\
5 & 8846.49 & 9479.29 & 11046.4 & 11830.1 & 17492.5 & 17638.3 & 14481.8 \\
6 & 9201.59 & 9873.76 & 11411 & 12488 & 15585.3 & 18152.4 & 15358.2 \\
7 & 9293.07 & 10015.1 & 11685.1 & 13029.5 & 17033.7 & 19868.5 & 16751.3 \\
8 & 10080.1 & 10993.2 & 12891.2 & 14295.8 & 17321.4 & 20518.4 & 17159.1 \\
9 & 9985.18 & 10842.3 & 12815.2 & 14613.6 & 18222.5 & 22332.9 & 19188 \\
10 & 10968.9 & 12001.6 & 14474.9 & 16602.5 & 19960 & 24353.3 & 21345.3 \\
\hline
%\bf{2022} &  &  &  & & & & \\
\multicolumn{8}{| c |}{\textbf{2023: September}}\\
\hline
1 & 8057.65 & 8039.26 & 7748.32 & 7366.06 & 7277.01 & 7381.23 & 7381.57 \\
2 & 7726.42 & 7990.22 & 8786.58 & 8828.21 & 14163.6 & 12041.9 & 10029.1 \\
3 & 8247.64 & 8609.77 & 9618.74 & 10150.4 & 14170.5 & 14536.1 & 12265.5 \\
4 & 8313.97 & 8631 & 9791.5 & 9795.11 & 18703.7 & 14820.3 & 11736 \\
5 & 8759.38 & 9204.49 & 10345.5 & 11259.8 & 13905.7 & 16652 & 14282 \\
6 & 8851.15 & 9357.15 & 10724.4 & 11653 & 16282.7 & 17375 & 14521.9 \\
7 & 9289.65 & 9878.35 & 11350.2 & 12612.5 & 15641.1 & 18960.8 & 16187 \\
8 & 9656.35 & 10336.8 & 12032.2 & 13588.6 & 16911.5 & 20589.9 & 17615.2 \\
9 & 10122.7 & 10879.8 & 12859.4 & 14816 & 18280.5 & 22704 & 19758.4 \\
10 & 11172.9 & 12165.6 & 14822.7 & 17351.5 & 20822.6 & 25700.6 & 22812.4 \\
\hline

%\bf{2023} &  &  &  & & & & \\
\multicolumn{8}{| c |}{\textbf{2023: November}}\\
\hline
1 & 7982.57 & 7976.11 & 7746.84 & 7371.42 & 7374.01 & 7416.82 & 7387.47 \\
2 & 7606.96 & 7848.7 & 8421.42 & 8450.72 & 11883.1 & 10576.5 & 9345.26 \\
3 & 8009.92 & 8271.76 & 9076.5 & 9113.51 & 14879.4 & 13007.5 & 10800.5 \\
4 & 8443.29 & 8714.75 & 9430.96 & 9844.52 & 11725.6 & 13884.5 & 12143.8 \\
5 & 8688.21 & 9043.79 & 10003.6 & 10602.4 & 13637.5 & 15634.6 & 13570.2 \\
6 & 8371.09 & 8671.7 & 9834.34 & 9692.7 & 18571.4 & 14782.6 & 11792.6 \\
7 & 8804.19 & 9223.59 & 10496.4 & 11154 & 16783.6 & 17135.9 & 14444 \\
8 & 9160.99 & 9627.78 & 10860.9 & 11853.9 & 14513.2 & 17931.3 & 15848 \\
9 & 9542.27 & 10134.7 & 11817.5 & 13258.5 & 17309 & 20476.2 & 18053.6 \\
10 & 11042.5 & 12012.4 & 14554.4 & 16720.1 & 20561.1 & 25084.9 & 22211.2 \\
\hline
\end{longtable}

Visualization of Earth landscapes using cartographic methods has many approaches, one of them is the use of \hl{RS} data. This method uses satellite images as input data for their processing and classification. Typically, cartography relies on the use of traditional GIS methods that can produce a visualization using a defined workflow dependent on a particular software. The detection of land cover types depicted on these maps is perceptible to the human eye depending on the spatial and temporal resolution of the \hl{RS} data. However, the use of automation and \hl{ML} approaches facilitates image processes, thanks to a more objective and accurate classification of satellite images.

\begin{figure}[H]
    \begin{center}
	\hspace{-35pt}\includegraphics[width=14.0cm]{Fig_09}
    \end{center}
    \caption{Accuracy evaluation using computed correctly classified pixels (\%), GRASS GIS.}
    \label{fig09}
\end{figure}

The maps describing land use in the Rif mountain region are scientific and decision-making tools for Moroccan authorities and educational institutions. The images processed by GRASS GIS methods can be used in research work as input data for modeling systems of bimonthly vegetation cycles in the geomorphological context of the Rif. In addition, it can also be used for operational monitoring of bimonthly vegetation changes in northern Morocco, to apply research recommendations that require precise knowledge of the Rif territories.

\begin{table}[H]
\centering
\begin{tabularx}{\textwidth}{| p{1.7cm} | X | X | X | X | X | X | }
  \toprule
   Iteration & 01/2023 & 03/2023 & 05/2023 & 07/2023 & 09/2023 & 11/2023  \\
   \midrule
    Iteration 1 & 74.15\% & 81.46\% & 58.44\% & 74.15\% & 81.12\% & 81.47\% \\
    Iteration 2 & 77.52\% & 88.60\% & 74.76\% & 77.52\% & 88.74\% & 78.51\% \\
    Iteration 3 & 90.02\% & 88.06\% & 89.85\% & 90.02\% & 92.05\% & 85.56\% \\
    Iteration 4 & 93.59\% & 93.42\% & 94.81\% & 93.59\% & 94.67\% & 92.87\% \\
    Iteration 5 & 95.31\% & 95.40\% & 96.52\% & 95.31\% & 95.25\% & 96.37\% \\
    Iteration 6 & 96.32\% & 95.94\% & 96.85\% & 96.32\% & 96.21\% & 97.33\% \\
    Iteration 7 & 97.11\% & 96.88\% & 97.34\% & 97.11\% & 97.09\% & 97.66\% \\
    Iteration 8 & 97.05\% & 97.78\% & 97.60\% & 97.05\% & 97.19\% & 97.97\% \\
    Iteration 9 & 97.55\% & 98.16\% & 98.03\% & 97.55\% & 97.46\% & 98.24\% \\
    \midrule
    Convergence & 98.1\% & 98.2\% & 98.0\% & 98.1\% & 98.0\% & 98.2\% \\
    \bottomrule
\end{tabularx}
\caption{Computed points stability for iterations during classification and convergence of classes: January, March, May, July, September and November 2023}\label{tabA06}
\end{table}

\hl{Complex interactions between human activities, climate change and environmental dynamics results in high vulnerability of landscapes and ecosystems in Morocco. This is particularly important in northern Morocco, a region which is highly sensitive to the contrasting effects of climate variability ranging from arid Sahara to humid Mediterranean, and where current land cover types are affected by pastoralism. As well as highlighting current distribution of land cover types in norther Morocco, using RS data processed by ML algorithms presents new maps that allow quantification of land cover changes computed in pixels. For instance, numerical analysis has revealed that even for the least modified areas of bare areas located near desert, regional inventories show slight changes over diverse seasonal periods during calendar year. Spatial biases are also apparent if comparing such land cover types as grassland, croplands and forests.}

%----------- section ----------->
\section{Conclusion}

The general objective of this paper aims to improve the production of land cover maps in northern Morocco \hl{using ML methods of }GRASS GIS. \hl{Our results clearly show that the land cover types in northern Morocco are changing. Nevertheless, the dynamics of such changes is vastly underrepresented in existing case studing on environmental mapping in norther Africa. At the same time, satellite images present the largest provider of information on the distribution of landscape patches and vegetation. To this end, we applied advanced }algorithms 'Decision trees' and 'Extra trees') \hl{for processing a} time series of \hl{the }Landsat 8-9 OLI/TIRS satellite images. 

\hl{More than this, however, we have also shown the technical difference between the} unsupervised and supervised classification (clustering) of \hl{RS} data. \hl{This is useful for further development of cartographic methods of RS data processing} in the paradigm of information processing systems\hl{. }Therefore, the specific objective of this paper aims to improve land cover mapping in northern Morocco from time series of Landsat satellite images using \hl{ML} methods by GRASS GIS and cartographic scripts. Two challenges have been identified in this work. The first concerns the technical choice of the classification algorithm (supervised “DecisionTreeClassifier” and “ExtraTreesClassifier” vs unsupervised clustering by MaxLike) using GRASS GIS, while the second concerns the \hl{identification} of seasonal \hl{bimonthly} landscape dynamics in the Rif Mountains of Morocco using\hl{ repetitive }data \hl{cycles} from January to November 2023\hl{. Comparison of these data enabled to reveal }the relationship between \hl{the growth of }vegetation and mountain geomorphology.

\hl{Taken together, this paper shows that the least well recorded region of the northern Morocco is largely affected by environmental changes: the comparison of presented maps showed a dynamics of agricultural and forest areas in the order of 7-10\% within the year of 2023, and a slight increase of bare soil of 8\% during the same year. Moreover, a notable change in the areas occupied by deciduous and mixed forest is detected. Currently, there are no land cover maps available for Morocco that depict bimonthly dynamics of land cover types for the region of Rif Mountains. In this regard, this paper contributes to improve this gap through presenting new maps for further evaluating of biophysical and environmental parameters in Northern Morocco. }

This paper demonstrated the implementation of the \hl{ML cartographic} approach to \hl{RS data processing, technically performed using GRASS GIS tools. The presented work enabled to }build a\hl{ }classification workflow of six multispectral satellite images covering\hl{ Northern} Morocco, \hl{around }Rif Mountains. \hl{Several} applications of the presented results and the GRASS GIS satellite image processing methods \hl{are }described in this paper. The \hl{demonstrated maps show dynamics in the }land cover in the northern\hl{ }Morocco \hl{which }presents the \hl{environmental dynamics that reflects the impact factors from regional }ecological, vegetal and geological \hl{setting in North Africa}.

%----------- section ----------->
%\authorcontributions{Supervision, conceptualization, methodology, software, resources, funding acquisition, and project administration, writing---original draft preparation, methodology, software, data curation, visualization, formal analysis, validation, writing---review and editing, and investigation, P.L.}

%\funding{The publication was funded by the Editorial Office of Geomatics, Multidisciplinary Digital Publishing Institute (MDPI), by providing 100\% discount for the APC of this manuscript.}

\institutionalreview{Not applicable.}

\informedconsent{Not applicable.}

\dataavailability{Not applicable} 

%\acknowledgments{The authors thank the reviewers for reading and review of this manuscript.}

\conflictsofinterest{The authors declare no conflict of interest.} 

\abbreviations{Abbreviations}{
The following abbreviations are used in this manuscript:\\

\noindent 
\begin{tabular}{@{}ll}
DTC & Decision Tree Classifier \\
ETC & Extra Trees Classifier \\
\hl{EO} & \hl{Earth Observation} \\
\hl{FAO} & \hl{Food and Agriculture Organization} \\
GRASS & Geographic Resources Analysis Support System \\
GIS & Geographic Information System \\
Landsat OLI/TIRS & Landsat Operational Land Imager and Thermal Infrared Sensor  \\
GEBCO & General Bathymetric Chart of the Oceans \\
GMT & Generic Mapping Tools \\
QGIS & Quantum Geographic Information System \\
ML & Machine Learning \\
RS & Remote Sensing \\
SWIR & Shortwave Infrared \\
NIR & Near Infrared \\
USGS & United States Geological Survey  \\
WGS84 & World Geodetic System 84 \\
\end{tabular}
}

%%%%%%%%%%%%%%%%%%%%%%%%%%%%%%%%%%%%%%%%%%
%% Optional
\appendixtitles{no} % Leave argument "no" if all appendix headings stay EMPTY (then no dot is printed after "Appendix A"). If the appendix sections contain a heading then change the argument to "yes".
\appendixstart
\appendix
%--------------------------->
\section[\appendixname~\thesection]{Initial means for classes of pixels in each multispectral band of the Landsat OLI/TIRS satellite images}

\begin{table}[H]
\centering
\begin{tabularx}{\textwidth}{| p{1.5cm} | X | X | X | X | X | X | X | }
  \toprule
   Class & Band 1 & Band 2 & Band 3 & Band 4 & Band 5 & Band 6 & Band 7  \\
   \midrule
    Class 1 & 7353.24 & 7533.21 & 8177.38 & 7818.66 & 10987.1 & 10158.3 & 8739.9 \\
    Class 2 & 7569.99 & 7774.52 & 8500.84 & 8237.08 & 11888.8 & 10959.1 & 9396.37 \\
    Class 3 & 7786.75 & 8015.83 & 8824.3 & 8655.49 & 12790.5 & 11759.9 & 10052.8 \\
    Class 4 & 8003.51 & 8257.14 & 9147.76 & 9073.91 & 13692.2 & 12560.7 & 10709.3 \\
    Class 5 & 8220.27 & 8498.45 & 9471.21 & 9492.32 & 14593.9 & 13361.5 & 11365.8 \\
    Class 6 & 8437.03 & 8739.76 & 9794.67 & 9910.73 & 15495.6 & 14162.3 & 12022.2 \\
    Class 7 & 8653.78 & 8981.07 & 10118.1 & 10329.1 & 16397.3 & 14963.1 & 12678.7 \\
    Class 8 & 8870.54 & 9222.38 & 10441.6 & 10747.6 & 17299 & 15763.9 & 13335.2 \\
    Class 9 & 9087.3 & 9463.69 & 10765 & 11166 & 18200.8 & 16564.7 & 13991.6 \\
    Class 10 & 9304.06 & 9705 & 11088.5 & 11584.4 & 19102.5 & 17365.5 & 14648.1 \\
    \bottomrule
\end{tabularx}
\caption{January 2023: initial means for each multispectral band of Landsat OLI/TIRS images}\label{tabA01}
\end{table}

\begin{table}[H]
\centering
\begin{tabularx}{\textwidth}{| p{1.5cm} | X | X | X | X | X | X | X | }
  \toprule
   Class & Band 1 & Band 2 & Band 3 & Band 4 & Band 5 & Band 6 & Band 7  \\
   \midrule
    Class 1 & 8036.12 & 8397.21 & 9295.02 & 9561.14 & 13138.1 & 12962.9 & 10897.1 \\
    Class 2 & 8246.35 & 8665.51 & 9702.28 & 10172.8 & 14019.7 & 13972.2 & 11692.2 \\
    Class 3 & 8456.57 & 8933.81 & 10109.5 & 10784.5 & 14901.3 & 14981.5 & 12487.3 \\
    Class 4 & 8666.8 & 9202.11 & 10516.8 & 11396.2 & 15783 & 15990.8 & 13282.5 \\
    Class 5 & 8877.03 & 9470.41 & 10924.1 & 12007.9 & 16664.6 & 17000.2 & 14077.6 \\
    Class 6 & 9087.26 & 9738.71 & 11331.3 & 12619.6 & 17546.2 & 18009.5 & 14872.7 \\
    Class 7 & 9297.49 & 10007 & 11738.6 & 13231.3 & 18427.9 & 19018.8 & 15667.8 \\
    Class 8 & 9507.72 & 10275.3 & 12145.8 & 13842.9 & 19309.5 & 20028.1 & 16463 \\
    Class 9 & 9717.94 & 10543.6 & 12553.1 & 14454.6 & 20191.1 & 21037.4 & 17258.1 \\
    Class 10 & 9928.17 & 10811.9 & 12960.4 & 15066.3 & 21072.7 & 22046.7 & 18053.2 \\
    \bottomrule
\end{tabularx}
\caption{March 2023: initial means for each multispectral band of Landsat OLI/TIRS satellite images}\label{tabA02}
\end{table}

\begin{table}[H]
\centering
\begin{tabularx}{\textwidth}{| p{1.5cm} | X | X | X | X | X | X | X | }
  \toprule
   Class & Band 1 & Band 2 & Band 3 & Band 4 & Band 5 & Band 6 & Band 7  \\
   \midrule
    Class 1 & 7598.78 & 7764.07 & 8434.53 & 7941.99 & 12295.6 & 10618.7 & 8849.61 \\
    Class 2 & 7754.07 & 7949.13 & 8718.98 & 8335.56 & 13270.4 & 11401 & 9511.46 \\
    Class 3 & 7909.36 & 8134.19 & 9003.43 & 8729.14 & 14245.2 & 12183.3 & 10173.3 \\
    Class 4 & 8064.65 & 8319.25 & 9287.89 & 9122.72 & 15219.9 & 12965.7 & 10835.2 \\
    Class 5 & 8219.95 & 8504.31 & 9572.34 & 9516.29 & 16194.7 & 13748 & 11497 \\
    Class 6 & 8375.24 & 8689.37 & 9856.79 & 9909.87 & 17169.5 & 14530.3 & 12158.8 \\
    Class 7 & 8530.53 & 8874.43 & 10141.2 & 10303.4 & 18144.3 & 15312.6 & 12820.7 \\
    Class 8 & 8685.82 & 9059.49 & 10425.7 & 10697 & 19119.1 & 16095 & 13482.5 \\
    Class 9 & 8841.12 & 9244.55 & 10710.1 & 11090.6 & 20093.9 & 16877.3 & 14144.4 \\
    Class 10 & 8996.41 & 9429.61 & 10994.6 & 11484.2 & 21068.7 & 17659.6 & 14806.2 \\
    \bottomrule
\end{tabularx}
\caption{May 2023: initial means for each multispectral band of Landsat OLI/TIRS satellite images}\label{tabA03}
\end{table}

\begin{table}[H]
\centering
\begin{tabularx}{\textwidth}{| p{1.5cm} | X | X | X | X | X | X | X | }
  \toprule
   Class & Band1 & Band 2 & Band 3 & Band 4 & Band 5 & Band 6 & Band 7  \\
   \midrule
    Class 1 & 8146.96 & 8572.74 & 9528.06 & 9871.64 & 12802.2 & 13459.3 & 11517.8 \\
    Class 2 & 8370.6 & 8848.08 & 9926.18 & 10412.8 & 13530.9 & 14431.9 & 12322.4 \\
    Class 3 & 8594.23 & 9123.41 & 10324.3 & 10954 & 14259.6 & 15404.5 & 13127.1 \\
    Class 4 & 8817.87 & 9398.75 & 10722.4 & 11495.2 & 14988.3 & 16377.2 & 13931.7 \\
    Class 5 & 9041.5 & 9674.09 & 11120.5 & 12036.4 & 15717.1 & 17349.8 & 14736.3 \\
    Class 6 & 9265.14 & 9949.43 & 11518.6 & 12577.6 & 16445.8 & 18322.4 & 15541 \\
    Class 7 & 9488.77 & 10224.8 & 11916.7 & 13118.9 & 17174.5 & 19295.1 & 16345.6 \\
    Class 8 & 9712.41 & 10500.1 & 12314.9 & 13660.1 & 17903.2 & 20267.7 & 17150.3 \\
    Class 9 & 9936.04 & 10775.4 & 12713 & 14201.3 & 18631.9 & 21240.3 & 17954.9 \\
    Class 10 & 10159.7 & 11050.8 & 13111.1 & 14742.5 & 19360.6 & 22212.9 & 18759.6 \\
    \bottomrule
\end{tabularx}
\caption{July 2023: initial means for each multispectral band of Landsat OLI/TIRS satellite images}\label{tabA04}
\end{table}

\begin{table}[H]
\centering
\begin{tabularx}{\textwidth}{| p{1.5cm} | X | X | X | X | X | X | X | }
  \toprule
   Class & Band1 & Band 2 & Band 3 & Band 4 & Band 5 & Band 6 & Band 7  \\
   \midrule
    Class 1 & 8037.45 & 8299.79 & 8964.58 & 9252.55 & 11892.9 & 12555.5 & 10819.7 \\
    Class 2 & 8252.53 & 8565.39 & 9363.68 & 9802.02 & 12631.3 & 13569.5 & 11680.2 \\
    Class 3 & 8467.62 & 8830.99 & 9762.77 & 10351.5 & 13369.8 & 14583.6 & 12540.8 \\
    Class 4 & 8682.7 & 9096.59 & 10161.9 & 10901 & 14108.2 & 15597.6 & 13401.3 \\
    Class 5 & 8897.79 & 9362.19 & 10561 & 11450.4 & 14846.6 & 16611.6 & 14261.9 \\
    Class 6 & 9112.87 & 9627.78 & 10960 & 11999.9 & 15585.1 & 17625.6 & 15122.4 \\
    Class 7 & 9327.96 & 9893.38 & 11359.1 & 12549.4 & 16323.5 & 18639.6 & 15983 \\
    Class 8 & 9543.04 & 10159 & 11758.2 & 13098.9 & 17061.9 & 19653.6 & 16843.5 \\
    Class 9 & 9758.13 & 10424.6 & 12157.3 & 13648.3 & 17800.4 & 20667.6 & 17704.1 \\
    Class 10 & 9973.21 & 10690.2 & 12556.4 & 14197.8 & 18538.8 & 21681.7 & 18564.6 \\
    \bottomrule
\end{tabularx}
\caption{September 2023: initial means for each multispectral band of Landsat OLI/TIRS images}\label{tabA05}
\end{table}

\begin{table}[H]
\centering
\begin{tabularx}{\textwidth}{| p{1.5cm} | X | X | X | X | X | X | X | }
  \toprule
   Class & Band 1 & Band 2 & Band 3 & Band 4 & Band 5 & Band 6 & Band 7  \\
   \midrule
    Class 1 & 7805.3 & 7978.69 & 8394.85 & 8339.26 & 10618.8 & 10876.8 & 9549.69 \\
    Class 2 & 7994.42 & 8207.25 & 8745.73 & 8818 & 11398.4 & 11816.7 & 10351.8 \\
    Class 3 & 8183.55 & 8435.81 & 9096.6 & 9296.74 & 12178 & 12756.6 & 11154 \\
    Class 4 & 8372.67 & 8664.38 & 9447.48 & 9775.48 & 12957.6 & 13696.6 & 11956.1 \\
    Class 5 & 8561.79 & 8892.94 & 9798.36 & 10254.2 & 13737.2 & 14636.5 & 12758.2 \\
    Class 6 & 8750.92 & 9121.5 & 10149.2 & 10733 & 14516.8 & 15576.5 & 13560.4 \\
    Class 7 & 8940.04 & 9350.07 & 10500.1 & 11211.7 & 15296.4 & 16516.4 & 14362.5 \\
    Class 8 & 9129.17 & 9578.63 & 10851 & 11690.4 & 16076 & 17456.4 & 15164.7 \\
    Class 9 & 9318.29 & 9807.2 & 11201.9 & 12169.2 & 16855.6 & 18396.3 & 15966.8 \\
    Class 10 & 9507.41 & 10035.8 & 11552.8 & 12647.9 & 17635.2 & 19336.3 & 16768.9 \\
    \bottomrule
\end{tabularx}
\caption{November 2023: initial means for each multispectral band of Landsat OLI/TIRS images}\label{tabA06}
\end{table}

%--------------------------->
\section[\appendixname~\thesection]{Standard deviations for DNs of pixels by 10 computed classes}

\begin{longtable}{@{\extracolsep{\fill}}p{1.5cm}cccccccc@{}}
\caption{Bimonthly evaluation of the standard deviations of Digital Numbers (DN) computed for pixels assigned to 10 land cover classes in landscapes of Rif Mountains, north Morocco. Evaluated each multispectral band (1 to 7) of Landsat 8-9 OLI/TIRS. Period: 2023, January to May.} \label{tabA2}\\
\hline \textbf{Class} & \textbf{Band 1} & \textbf{Band 2} & \textbf{Band 3} & \textbf{Band 4} & \textbf{Band 5} & \textbf{Band 6} & \textbf{Band 7} \\ \hline 
\endfirsthead

\multicolumn{5}{c}%
{{\tablename\ \thetable{} -- continued from previous page}} \\
\textbf{Class} & \textbf{Band 1} & \textbf{Band 2} & \textbf{Band 3} & \textbf{Band 4} & \textbf{Band 5} & \textbf{Band 6} & \textbf{Band 7}  \\ \hline 
\endhead

\hline \multicolumn{8}{r}{{Continued on next page}} \\ \hline
\endfoot

\hline \hline
\endlastfoot

%\bf{2019} &  &  &  &  &  &  & \\
\multicolumn{8}{| c |}{\textbf{2023: January}}\\
\hline
1 & 520.366 & 535.906 & 762.029 & 534.589 & 458.971 & 261.960 & 185.243 \\
2 & 490.239 & 471.037 & 624.903 & 718.469 & 1108.75 & 864.994 & 643.283 \\
3 & 469.625 & 475.629 & 523.194 & 551.304 & 924.17 & 813.025 & 684.606 \\
4 & 314.352 & 297.355 & 402.396 & 473.559 & 946.247 & 881.536 & 697.401 \\
5 & 355.762 & 363.823 & 482.001 & 628.670 & 1123.910 & 809.348 & 943.121 \\
6 & 398.402 & 397.373 & 474.544 & 609.048 & 1059.400 & 1019.420 & 849.959 \\
7 & 831.164 & 913.648 & 997.848 & 1243.020 & 1609.220 & 1416.340 & 1089.100 \\
8 & 602.747 & 649.726 & 694.899 & 843.584 & 1472.750 & 926.428 & 972.410 \\
9 & 563.373 & 645.449 & 846.794 & 1005.060 & 1844.240 & 1204.890 & 1309.140 \\
10 & 3912.350 & 4022.450 & 3523.310 & 3432.850 & 3098.910 & 2796.470 & 2599.350 \\
\hline
%\bf{2022} &  &  &  & & & & \\
\multicolumn{8}{| c |}{\textbf{2023: March}}\\
\hline
1 & 490.469 & 541.212 & 879.207 & 716.637 & 616.488 & 686.935 & 613.483 \\
2 & 255.312 & 246.682 & 318.072 & 462.128 & 1664.020 & 929.262 & 719.619 \\
3 & 541.560 & 602.780 & 686.249 & 782.835 & 1273.070 & 980.419 & 834.318 \\
4 & 344.733 & 387.283 & 524.843 & 812.884 & 1659.040 & 980.065 & 855.260 \\
5 & 377.657 & 422.529 & 565.366 & 724.443 & 1044.500 & 753.815 & 856.870 \\
6 & 507.082 & 559.630 & 652.069 & 974.570 & 1023.500 & 866.568 & 784.542 \\
7 & 449.736 & 488.873 & 597.974 & 692.265 & 1060.560 & 878.137 & 929.794 \\
8 & 532.728 & 589.376 & 727.215 & 948.085 & 1317.450 & 934.689 & 1002.260 \\
9 & 507.096 & 562.500 & 709.818 & 853.806 & 1051.740 & 1083.610 & 1336.400 \\
10 & 753.846 & 857.879 & 1039.350 & 1108.080 & 1502.720 & 1502.480 & 2189.920 \\
\hline

%\bf{2023} &  &  &  & & & & \\
\multicolumn{8}{| c |}{\textbf{2023: May}}\\
\hline
1 & 465.095 & 566.529 & 901.861 & 522.758 & 321.154 & 193.603 & 120.269 \\
2 & 277.062 & 255.083 & 376.608 & 450.477 & 1258.88 & 804.346 & 710.667 \\
3 & 441.646 & 468.557 & 562.157 & 597.613 & 1136.880 & 723.626 & 673.917 \\
4 & 378.651 & 385.161 & 519.410 & 649.726 & 1220.780 & 868.215 & 1039.880 \\
5 & 258.976 & 264.420 & 393.296 & 494.295 & 1055.540 & 951.268 & 799.838 \\
6 & 228.533 & 221.204 & 444.333 & 414.391 & 1722.940 & 994.491 & 630.409 \\
7 & 531.397 & 586.941 & 678.976 & 806.428 & 1337.73 & 1061.7 & 936.166 \\
8 & 623.942 & 671.415 & 760.778 & 817.799 & 1089.700 & 905.306 & 884.892 \\
9 & 528.555 & 602.147 & 767.459 & 860.058 & 1806.600 & 1088.830 & 1165.730 \\
10 & 1210.190 & 1398.230 & 1656.380 & 1682.400 & 1778.850 & 1867.540 & 2167.590 \\
\hline
\end{longtable}

\begin{longtable}{@{\extracolsep{\fill}}p{1.5cm}cccccccc@{}}
\caption{Bimonthly evaluation of the standard deviations of Digital Numbers (DN) computed for pixels assigned to 10 land cover classes in landscapes of Rif Mountains, north Morocco. Evaluated each multispectral band (1 to 7) of Landsat 8-9 OLI/TIRS. Period: 2023, July to November.} \label{tab04}\\
\hline \textbf{Class} & \textbf{Band 1} & \textbf{Band 2} & \textbf{Band 3} & \textbf{Band 4} & \textbf{Band 5} & \textbf{Band 6} & \textbf{Band 7} \\ \hline 
\endfirsthead

\multicolumn{5}{c}%
{{\tablename\ \thetable{} -- continued from previous page}} \\
\textbf{Class} & \textbf{Band 1} & \textbf{Band 2} & \textbf{Band 3} & \textbf{Band 4} & \textbf{Band 5} & \textbf{Band 6} & \textbf{Band 7}  \\ \hline 
\endhead

\hline \multicolumn{8}{r}{{Continued on next page}} \\ \hline
\endfoot

\hline \hline
\endlastfoot

%\bf{2019} &  &  &  &  &  &  & \\
\multicolumn{8}{| c |}{\textbf{2023: July}}\\
\hline
1 & 401.603 & 474.680 & 753.823 & 554.671 & 452.066 & 377.558 & 279.117 \\
2 & 444.544 & 494.063 & 605.472 & 689.234 & 1469.020 & 1074.000 & 858.682 \\
3 & 400.815 & 449.109 & 548.777 & 721.564 & 1773.570 & 1030.710 & 912.245 \\
4 & 329.822 & 337.588 & 456.271 & 617.552 & 934.712 & 737.397 & 693.313 \\
5 & 383.645 & 392.394 & 484.058 & 558.823 & 955.127 & 757.462 & 640.994 \\
6 & 369.823 & 422.794 & 574.353 & 625.855 & 807.736 & 699.93 & 634.601 \\
7 & 380.957 & 393.860 & 462.246 & 493.918 & 938.201 & 756.140 & 818.874 \\
8 & 389.170 & 433.991 & 512.028 & 543.331 & 810.822 & 715.835 & 672.958 \\
9 & 585.918 & 680.160 & 803.684 & 738.175 & 743.199 & 762.175 & 944.871 \\
10 & 994.099 & 1126.100 & 1284.760 & 1153.090 & 1115.620 & 1242.490 & 1575.490 \\
\hline
%\bf{2022} &  &  &  & & & & \\
\multicolumn{8}{| c |}{\textbf{2023: September}}\\
\hline
1 & 437.651 & 443.108 & 680.601 & 464.861 & 251.267 & 134.622 & 109.822 \\
2 & 352.790 & 343.761 & 458.110 & 604.993 & 1647.410 & 1009.150 & 737.153 \\
3 & 413.844 & 428.298 & 527.591 & 627.276 & 1270.460 & 820.868 & 669.725 \\
4 & 609.427 & 665.603 & 750.628 & 922.569 & 2049.290 & 1237.420 & 1107.210 \\
5 & 410.913 & 444.970 & 557.272 & 605.301 & 936.943 & 802.527 & 790.452 \\
6 & 412.096 & 463.575 & 605.626 & 711.301 & 1026.750 & 826.580 & 710.861 \\
7 & 476.481 & 556.664 & 690.121 & 675.737 & 961.119 & 727.211 & 750.357 \\
8 & 560.709 & 661.131 & 788.185 & 748.966 & 878.300 & 826.902 & 929.749 \\
9 & 734.715 & 872.060 & 1021.830 & 874.668 & 875.353 & 1032.820 & 1289.020 \\
10 & 1228.510 & 1477.600 & 1759.470 & 1582.050 & 1479.630 & 1594.770 & 1825.860 \\
\hline

%\bf{2023} &  &  &  & & & & \\
\multicolumn{8}{| c |}{\textbf{2023: November}}\\
\hline
1 & 408.799 & 366.029 & 518.468 & 358.299 & 494.063 & 288.290 & 183.250 \\
2 & 337.779 & 313.223 & 378.857 & 471.889 & 1236.690 & 864.664 & 718.244 \\
3 & 329.228 & 309.538 & 391.118 & 526.252 & 1133.950 & 971.402 & 842.809 \\
4 & 357.573 & 340.276 & 405.652 & 474.368 & 868.293 & 878.208 & 719.913 \\
5 & 407.823 & 418.715 & 501.585 & 583.480 & 1177.420 & 746.789 & 760.614 \\
6 & 560.647 & 596.644 & 680.376 & 807.904 & 1923.630 & 1063.100 & 992.282 \\
7 & 404.268 & 434.357 & 567.739 & 689.356 & 1145.220 & 911.636 & 828.516 \\
8 & 502.947 & 558.014 & 686.296 & 677.680 & 976.150 & 885.420 & 857.641 \\
9 & 571.593 & 667.463 & 845.177 & 950.398 & 1288.210 & 1164.560 & 1261.840 \\
10 & 1130.670 & 1394.490 & 1762.720 & 1842.670 & 1953.210 & 2219.250 & 2058.000 \\
\hline
\end{longtable}

%%%%%%%%%%%%%%%%%%%%%%%%%%%%%%%%%%%%%%%%%%
\begin{adjustwidth}{-\extralength}{0cm}
%\printendnotes[custom] % Un-comment to print a list of endnotes

\reftitle{References}
%\nocite{*}

%=====================================
% References, variant A: external bibliography
%=====================================
\bibliography{Morocco.bib}

%%%%%%%%%%%%%%%%%%%%%%%%%%%%%%%%%%%%%%%%%%
\end{adjustwidth}
\end{document}
